\clearpage\SourceLeaf{199}%
\pdfbookmark[1]{Gauss an Schumacher, 23. November 1847}{schumacher18471123}
\fancyhead[L]{\small GAUSS AN SCHUMACHER}
\begin{center}
{\small GAUSS AN SCHUMACHER.\par}
\medskip\hrule\medskip
{\small Astronomische Nachrichten herausgegeben von H.\ C.\ \textsc{Schumacher}.\par
Band 27.\ Nr.\ 625.\ S.\ 1.\par}
\medskip\hrule
\end{center}\medskip
\begin{flushright}Göttingen, 1847 November 23.\end{flushright}
Da die in Nr.\ 615 der A.\ N.\ abgedruckten zweiten Elemente der Iris, welche Herr Prof.\ \textsc{Goldschmidt} gleich nach dem 19.\ September berechnet hatte, auch noch in den folgenden Monaten eine gute Übereinstimmung mit den Beobachtungen gezeigt haben (am 17.\ und 18.\ November war die Differenz von Herrn \textsc{Rümcker}'s Meridianbeobachtung $16''$ in ger.\ Aufst.\ und $2''$ in der Abweichung), so ist man berechtigt, sie als schon sehr genähert zu betrachten, und ich habe deshalb Herrn Prof.\ \textsc{Goldschmidt} veranlasst, danach den Zodiacus der Iris zu berechnen. Das Resultat dieser Arbeit, welches zur Erleichterung der Nachforschungen auf Identität in frühern Beobachtungen wird dienen können, lasse ich hier nachfolgen.

Indem ich bei dieser Veranlassung den Aufsatz wieder in die Hände nehme, in welchem ich vor 44 Jahren die allgemeine Methode zur Bestimmung der Limiten eines solchen Zodiacus gegeben habe, sehe ich, dass ich darin auch schon die Ausnahmsfälle angedeutet habe, wo das Feld der geocentrischen Erscheinung eines die Sonne nach \textsc{Kepler}'schen Gesetzen umkreisenden Himmelskörpers auf der Himmelskugel entweder gar keine oder nur Eine Limite hat, obgleich die Methode im erstern Falle eine in sich zurücklaufende Linie, im andern zwei solche Linien ergibt. Auch ist die Frage daselbst aufgeworfen, was denn in solchen Ausnahmsfällen diese durch Rechnung gefundenen Linien eigentlich
\SourceLeaf{200}%
bedeuten. Ich habe mich damals auf diese Andeutungen beschränkt, weil eine weitere Ausführung dort ein Horsd'oeuvre gewesen wäre, und ich auch gern andern das Vergnügen lassen wollte, sich mit einer meiner Meinung nach nicht uninteressanten mathematischen Aufgabe zu beschäftigen. Da mir jedoch nicht bekannt geworden ist, dass ein anderer in der langen Zwischenzeit die mir inzwischen ganz aus dem Gedächtniss gekommene Frage aufgenommen hätte, so ergreife ich diese Gelegenheit, um wenigstens den Hauptnerv des zur Beantwortung nöthigen hier mitzutheilen.

Ein geocentrischer Ort des in Rede stehenden Planeten (oder Cometen) geht hervor, indem man einen Punkt der Bahn des letztern mit einem Punkt der Erdbahn combinirt; es kann aber auch einerlei geocentrischer Ort aus zwei, drei, oder vier verschiedenen Combinationen hervorgehen. Um die Vorstellungen zu fixiren, lege man durch die Sonne eine Ebene, gegen welche die einen vorgegebenen geocentrischen Ort, auf der Himmelskugel, repräsentirende gerade Linie normal ist, und projicire orthographisch auf diese Ebene sowohl die Erdbahn als die Planetenbahn. Beide Projectionen sind Ellipsen, oder allgemeiner Kegelschnitte, die sich entweder gar nicht, oder zweimal oder viermal schneiden werden; eine Berührung ist dabei wie das Verschmelzen zweier Schnitte zu betrachten. Durch jeden Schnitt wird ein geocentrischer Ort vorgestellt, der entweder mit dem vorgegebenen identisch, oder ihm auf der Himmelskugel entgegengesetzt ist, je nachdem von den beiden Punkten der Planetenbahn und der Erdbahn, deren Projection zusammenfällt, der erstere oder der andere höherliegend ist; als obere Seite der Projectionsebene diejenige betrachtet, auf welcher der vorgegebene geocentrische Ort liegt.

Es ist hieraus klar, dass wenn irgend ein Punkt der Himmelskugel. fragweise als geocentrischer Ort, aufgegeben wird, er dies entweder auf gar keine Weise oder auf Eine Art, oder auf zwei, drei oder vier Arten sein kann. Für einen gegebenen Planeten scheidet sich so die ganze Fläche der Himmelskugel in verschiedene Theile, und die nach der von mir gegebenen Methode bestimmten Linien sind, allgemein zu reden, nichts anders, als Scheidungen zwischen zwei Flächentheilen der Kugel, wo die auf der einen Seite liegenden Punkte auf zwei Arten mehr, geocentrische Örter sein können als die auf der andern. Die in den Scheidungslinien liegenden Punkte machen den Übergang, d.\ i.\ sie sind auf Eine Art mehr als die Punkte auf der einen Seite, und auf Eine Art weniger als die
\SourceLeaf{201}%
Punkte auf der andern Seite fähig geocentrische Örter zu sein. Übrigens lässt sich auch das Criterium angeben, wonach a priori entschieden wird, auf welcher Seite der Scheidungslinie zwei Auflösungen mehr Statt finden als auf der andern, wobei ich jedoch gegenwärtig mich nicht aufhalten will.

Alle bisher bekannten Planeten haben solche Bahnen, dass die durch die Theorie gefundenen Scheidungslinien immer wahre Limiten sind. Es sind nämlich zwei Scheidungslinien, die die Himmelskugel in drei Flächenräume abtheilen; zwei sind isolirte Flächen, in welche gar keine geocentrische Örter fallen, der dritte zwischen jenen, gürtelartig, enthält alle geocentrischen Örter, und jeder Punkt innerhalb des Gürtels kann auf zwei Arten, jeder Punkt auf der Limite auf Eine Art geocentrischer Ort sein.

Es lassen sich aber fingirte Bahnen denken (Cometenbahnen werden vielleicht mehrere in dem Fall sein, worüber ich jedoch bisher keine Nachforschungen gemacht habe), wo zwei Limiten die Himmelskugel auch in drei Stücke scheiden, und wo der eine Theil gar keine geocentrischen Örter enthält, der andere die geocentrischen zu zwei Arten, der dritte hingegen die Punkte die auf vier verschiedene Arten geocentrische Örter sein können.

Noch weitere Mannigfaltigkeit ergibt sich, wenn Eine Scheidungslinie sich selbst einmal oder zweimal schneidet, eine einfache oder doppelte Schleife bildet. Im ersten Fall wird sie zwei, im zweiten drei Flächenräume von dem gürtelförmigen Theile abtrennen, in denen resp.\ 4 und 0 oder 4, 0, 4 Auflösungsarten Statt finden.

Noch anders verhält es sich mit einer Bahn, die in die Erdbahn wie ein Kettenring eingreift. In einem solchen Fall ist nur Eine zusammenhängende Limitenlinie, also zwei geschiedene Flächentheile, wo die Punkte des einen Theils auf Eine Art, die des andern Theils auf drei Arten geocentrische Örter sein können.
\par\medskip\centerline{\rule{.25\textwidth}{.3pt}}
